\appendix
\section{Source correspondence and verification boundary}
\label{app:sources}
This manuscript expands the frozen proof
\pathref{research_math/vertex_excess_sharp_exponent_20261008/SHARP_VERTEX_EXCESS_EXPONENT.md}
and its geometric-sharpness addendum. The frozen main has SHA-256
\pathref{8da4c012d4099465b4702ce4091ba4d1e2cc7862c12e8eb55cbd7467414e0f05}.
The separate traditional review is
\pathref{research_math/vertex_excess_exponent_independent_audit_20261008/REFEREE_REPORT.md},
SHA-256
\pathref{d638e1ebf4505f480d99c08c21a2012fb36f0e021cd86d828d142effd5df283b}.
The supplied review checks the main exponent, the exact actual-defect
constant, the arbitrary-prescribed-maximum quantifier, both defect
endpoints, and the lower family. It is an independent model's analytic
review, not a human referee report or a new Lean replay.

\subsection{Retained arguments reproduced here}
The volume-weighted determinant selection is retained from
\pathref{notes/quadratic-dimensional-simplex-stability/imports/weighted_anchors.md}
in the project publication tree. Section~\ref{sec:assignment} supplies its
finite-support seminorm proof. The bounded-vertex retention argument is
retained from
\pathref{research_math/linear_retention_investigation_20261008/BOUNDED_VERTEX_EXCESS_RETENTION.md}.
Section~\ref{sec:retention} includes both that proof and the unrestricted
retention proof. The exact truncation calculation is retained from
\pathref{notes/simplex-truncation-stability/proof.tex} in the publication
tree; Section~\ref{sec:sharpness} includes the facet calculation, the
$m=2$ case, the maximum-simplex proof, and the pyramid identities.
The package's \pathref{SOURCE_MAP.md} records the precise local source
paths, hashes, and section-level correspondence, and includes frozen
copies of the main input documents.

\subsection{Lean references refer to retained statements only}
For the earlier formal source, the module filenames and theorem namespaces
must be distinguished. In particular, the retained interfaces include
\begin{itemize}[leftmargin=1.5em]
\item \pathref{Entry005.finite_halfspace_entryA_lifted_moment}, in
\pathref{Entry005/PyramidEntryDefect.lean};
\item \pathref{Entry005.finite_halfspace_defect_over_first_moment}, in
\pathref{Entry005/PyramidMomentDefect.lean};
\item \pathref{Entry005.unit_ball_seminorm_first_moment_assignment}, in
\pathref{Entry005/SeminormFirstMoment.lean}.
\end{itemize}
Actual finite-law probability, centering, ball support, and brightness
interfaces are in \pathref{Entry005/FiniteHalfspaceConeLaw.lean};
finite-support assignment interfaces are in
\pathref{Entry005/FiniteBodyWeightedAssignment.lean}; the earlier
unrestricted prescribed-simplex interface is in
\pathref{Entry005/RelativeSimplexRetention.lean}.
These are references to retained code, not new formalization claims for
this manuscript. In particular, \pathref{PyramidEntryDefect} and
\pathref{PyramidMomentDefect} are module names, not the theorem namespace.

The older compact-limit main theorem does not by itself provide the
actual-facet incidence used in Corollary~\ref{cor:caps}. Here that incidence
is proved directly by choosing anchors from the finite law
\eqref{eq:law}. The new cone/prismoid estimate, ray-collapse/counting step,
finite-law specialization with this incidence, and final exponent assembly
still need corresponding Lean formalization. An earlier successful
\pathref{compiledMain} or other main build cannot certify those additions.

\subsection{Editorial clarifications implemented}
The positive-integer restrictions in Lemma~\ref{lem:cone}, unit projection
convention and homogeneous brightness notation, explicit two-dimensional
core, actual Lean namespace references, and mixed-volume factorial
convention implement all nonblocking clarifications requested by the
supplied review. The finite actual-facet chain is expanded in full in
Section~\ref{sec:enclosure}. Standard mixed-volume inputs remain credited
as classical, and no strict-monotonicity theorem from~\cite{BS} is needed.

The theorem is a mathematical manuscript candidate. Neither the focused
literature check nor the separate model review establishes priority,
journal suitability, or acceptance by professional mathematicians.
